\chapter{\texorpdfstring{$\beta$}{Beta}-eliminatie}\label{ch:b1:elimination}

Behandel 2-broom-2-methylpropaan met natriumethoxide in ethanol bij
kamertemperatuur, en het hoofdproduct is een ether; verwarm hetzelfde
mengsel, en het hoofdproduct is een alkeen, 2-methylpropeen, zonder
ethoxygroep. Het ethoxide-ion heeft niet als \omterm{def:b1:electronic-effects:nucleophile}{nucleofiel} opgetreden, dat
een koolstofatoom aanvalt, maar als base, die een proton wegneemt van het
naburige koolstofatoom terwijl het bromide vertrok. Deze concurrentie
tussen substitutie en eliminatie loopt door de hele organische synthese.
Dit hoofdstuk beschrijft de twee eliminatiemechanismen, hun opmerkelijke
geometrische voorwaarde, de regel die voorspelt welk alkeen ontstaat, en
hoe je een reactie de ene of de andere kant op stuurt.

\begin{recall}
SN1 en SN2 (\cref{ch:b1:nucleophilic-substitution}); \omterm{def:b1:stereochemistry-in-depth:representations}{newmanprojecties},
\omterm{def:b1:stereochemistry-in-depth:isomers}{conformaties} van cyclohexaan en de descriptoren $Z$/$E$
(\cref{ch:b1:stereochemistry-in-depth}); \omterm{def:b1:electronic-effects:intermediates}{carbokationen} en basen
(\cref{ch:b1:electronic-effects}).
\end{recall}

\section{Het E2-mechanisme}

\begin{definition}[\texorpdfstring{$\beta$}{Beta}-eliminatie, E2 en E1]\label{def:b1:elimination:elimination}
Een \emph{$\beta$-eliminatie}\index{beta-eliminatie@$\beta$-eliminatie}
verwijdert een \omterm{def:b1:electronic-effects:nucleophile}{vertrekkende groep} Y van een koolstofatoom (het
$\alpha$-koolstofatoom) en een proton van een naburig koolstofatoom (een
$\beta$-koolstofatoom), waarbij een dubbele binding tussen beide ontstaat.
Bij het \emph{E2-mechanisme}\index{E2-mechanisme} neemt een base het
proton weg terwijl Y vertrekt, in één \omterm{def:b1:elementary-steps:elementary-step}{elementaire stap}; bij het
\emph{E1-mechanisme}\index{E1-mechanisme} vertrekt Y eerst, wat een
\omterm{def:b1:electronic-effects:intermediates}{carbokation} geeft, dat daarna een $\beta$-proton verliest.
\end{definition}

\begin{proposition}[Snelheidswet van E2]\label{prop:b1:elimination:e2-law}
Voor E2 is $v = k[\mathrm{RY}][\mathrm{B}]$, waarin B de base is.
\end{proposition}

\begin{proof}
Eén bimoleculaire \omterm{def:b1:elementary-steps:elementary-step}{elementaire stap}, waarbij \omterm{def:b1:nucleophilic-substitution:substitution}{substraat} en base betrokken
zijn: de massawerkingswet voor \omterm{def:b1:elementary-steps:elementary-step}{elementaire stappen}
(\cref{ch:b1:elementary-steps}).
\end{proof}

\begin{definition}[Anti- en synperiplanair]\label{def:b1:elimination:periplanar}
Twee bindingen op naburige atomen zijn
\emph{antiperiplanair}\index{antiperiplanair} wanneer ze in één vlak aan
tegengestelde kanten liggen (\omterm{def:b1:stereochemistry-in-depth:conformations}{torsiehoek} $180^\circ$), en
\emph{synperiplanair}\index{synperiplanair} wanneer ze in één vlak aan
dezelfde kant liggen ($0^\circ$).
\end{definition}

\begin{omfigure}
\begin{tikzpicture}[font=\small, >={Stealth}]
  \pic (n) at (0,0) {omnewman={90}{270}};
  \node at (n-f1) {H}; \node at (n-f2) {\ce{CH3}}; \node at (n-f3) {H};
  \node at (n-b1) {Br}; \node at (n-b2) {\ce{CH3}}; \node at (n-b3) {H};
  \node (b) at (-2.6,2.0) {\ce{EtO-}};
  \draw[->, thick, omThm] (b) to[bend left=20] ($(n-f1)+(-0.15,0.1)$);
  \draw[->, thick, omThm] (0.12,0.8) to[bend left=40] (0.12,0.25);
  \draw[->, thick, omThm] (0.12,-0.6) to[bend right=35] ($(n-b1)+(0.25,0.15)$);
  \node[below] at (0,-2.0) {H en Br antiperiplanair};
  \draw[->, thick] (2.4,0) -- (3.6,0) node[midway, above] {E2};
  \node at (5.6,0) {\chemfig{C(-[:150]H_3C)(-[:210]H)=C(-[:30]H)(-[:-30]CH_3)}};
  \node at (5.6,-0.9) {\cip{E}-but-2-een};
\end{tikzpicture}

{\small E2-reactie van 2-broombutaan, bekeken langs de binding C3--C2 (C3
vooraan, C2 met Br achteraan) in een van haar twee \omterm{def:b1:elimination:periplanar}{antiperiplanaire}
\omterm{def:b1:stereochemistry-in-depth:isomers}{conformaties}; de andere, met het andere waterstofatoom van C3 anti ten
opzichte van Br, geeft \cip{Z}-but-2-een. Drie elektronenparen bewegen
samen: de base neemt het proton weg, het C--H-paar wordt de
$\pi$-binding, het C--Br-paar vertrekt met het broom. De groepen die in de
\omterm{def:b1:stereochemistry-in-depth:representations}{newmanprojectie} tegenover elkaar staan, eindigen aan dezelfde kant van de
dubbele binding.}
\end{omfigure}

\begin{proposition}[E2 is anti en stereospecifiek]\label{prop:b1:elimination:anti}
E2 verloopt vanuit de \omterm{def:b1:stereochemistry-in-depth:isomers}{conformatie} waarin de C--H- en de C--Y-binding
\omterm{def:b1:elimination:periplanar}{antiperiplanair} zijn. Ze is stereospecifiek: diastereomere \omterm{def:b1:nucleophilic-substitution:substitution}{substraten}
geven verschillende stereo-isomere alkenen, vastgelegd door de
anti-schikking.
\end{proposition}

\begin{proof}
In de \omterm{def:b1:elementary-steps:energy-profile}{overgangstoestand} moet het bindende C--H-paar het antibindende
gebied van de C--Y-binding binnengaan, terwijl de twee koolstofatomen
trigonaal worden; de overlapping is het best wanneer de twee bindingen
evenwijdig in één vlak liggen. De anti-schikking is gestaffeld (lage
energie) en houdt de base ver van de \omterm{def:b1:electronic-effects:nucleophile}{vertrekkende groep}; de syn-schikking
is eclips. Met H en Y anti liggen de twee andere groepen op elk
koolstofatoom vast: de groepen die in de \omterm{def:b1:stereochemistry-in-depth:representations}{newmanprojectie} aan dezelfde
kant staan, eindigen \textit{cis} op de dubbele binding. Een \omterm{def:b1:stereochemistry-in-depth:enantiomers}{diastereomeer}
heeft twee groepen op één koolstofatoom verwisseld, en geeft het andere
alkeen.
\end{proof}

\begin{proposition}[E2 op een cyclohexaan]\label{prop:b1:elimination:cyclohexane-e2}
Op een cyclohexaanring vereist E2 dat de \omterm{def:b1:electronic-effects:nucleophile}{vertrekkende groep} en het
$\beta$-waterstofatoom allebei \omterm{def:b1:stereochemistry-in-depth:conformations}{axiaal} staan (\textit{trans}-diaxiaal); een
\omterm{def:b1:electronic-effects:nucleophile}{vertrekkende groep} die \omterm{def:b1:stereochemistry-in-depth:conformations}{equatoriaal} vastgehouden wordt, kan pas reageren
als de ring omklapt.
\end{proposition}

\begin{proof}
Op naburige koolstofatomen van een \omterm{def:b1:stereochemistry-in-depth:conformations}{stoel} zijn twee \omterm{def:b1:stereochemistry-in-depth:conformations}{axiale} bindingen
\omterm{def:b1:elimination:periplanar}{antiperiplanair} (één omhoog, één omlaag, evenwijdig aan de as); een
\omterm{def:b1:stereochemistry-in-depth:conformations}{equatoriale} binding staat alleen anti ten opzichte van C--C-bindingen van
de ring, nooit ten opzichte van een C--H-binding.
\end{proof}

\begin{omfigure}
\begin{tikzpicture}[font=\small]
  \pic (a) at (0,0) {omchair};
  \draw[thick, omThm] (a-c1) -- (a-a1) node[above] {Cl};
  \draw[thick] (a-c1) -- (a-e1) node[right] {H};
  \draw[thick, omDef] (a-c2) -- (a-a2) node[below] {H};
  \draw[thick, omDef] (a-c6) -- (a-a6) node[below] {H};
  \node[align=left, anchor=west] at (3.0,0.2) {Cl axiaal (omhoog) op C1;\\axiale H (omlaag) op beide buren:\\twee antiperiplanaire H};
\end{tikzpicture}

{\small \textit{trans}-Diaxiale schikking op een \omterm{def:b1:stereochemistry-in-depth:conformations}{stoel}: het \omterm{def:b1:stereochemistry-in-depth:conformations}{axiale}
chlooratoom staat \omterm{def:b1:elimination:periplanar}{antiperiplanair} ten opzichte van de \omterm{def:b1:stereochemistry-in-depth:conformations}{axiale}
waterstofatomen van de twee naburige koolstofatomen (blauw).}
\end{omfigure}

\section{Het E1-mechanisme}

\begin{proposition}[Snelheidswet van E1]\label{prop:b1:elimination:e1-law}
Voor E1 is $v = k[\mathrm{RY}]$, waarbij de snelheid die van de ionisatie
is.
\end{proposition}

\begin{proof}
Zoals bij SN1 (\cref{prop:b1:nucleophilic-substitution:sn1-law}): het
\omterm{def:b1:electronic-effects:intermediates}{carbokation} wordt in een trage stap gevormd en snel verbruikt, hier door
verlies van een proton aan het oplosmiddel of aan een \omterm{def:b1:predominant-reaction:strong-acid}{zwakke base}; de
\omterm{def:b1:elementary-steps:qssa}{stationaire-toestandsbenadering} laat $v = k_1[\mathrm{RY}]$ over.
\end{proof}

E1 deelt zijn \omterm{def:b1:electronic-effects:intermediates}{carbokation} met SN1: tertiaire \omterm{def:b1:nucleophilic-substitution:substitution}{substraten} in \omterm{def:b1:intermolecular-forces-solvents:solvent-classes}{protische
oplosmiddelen} geven beide, waarbij het alkeen door verwarming begunstigd
wordt. Omdat het \omterm{def:b1:electronic-effects:intermediates}{carbokation} vlak is en het proton daarna verloren gaat,
heeft E1 geen geometrische voorwaarde en is het niet stereospecifiek; het
geeft vooral het stabielste alkeen, meestal het \cip{E}-alkeen.

\begin{omfigure}
\setchemfig{atom sep=2.1em}
\schemestart
\chemfig{C(-[2]CH_3)(-[4]H_3C)(-[6]CH_3)-Br}
\arrow{->[traag]}
\chemfig{C^{+}(-[2]CH_3)(-[:210]H_3C)(-[:-30]CH_3)}
\arrow{->[snel][\ce{-H+}]}
\chemfig{C(-[2]CH_3)(-[:210]H_3C)=[:-30]CH_2}
\schemestop
\par\vspace{4pt}

{\small E1: ionisatie van 2-broom-2-methylpropaan tot het \omterm{def:b1:electronic-effects:intermediates}{carbokation},
daarna verlies van een proton van een methylgroep aan het oplosmiddel,
wat 2-methylpropeen geeft.}
\end{omfigure}

\section{Regioselectiviteit: de regel van Zaitsev}

\begin{definition}[Regioselectieve en stereoselectieve reacties]\label{def:b1:elimination:selectivity}
Een reactie die meerdere \omterm{def:b1:stereochemistry-in-depth:isomers}{constitutie-isomeren} kan geven maar er vooral
één vormt, is een
\emph{regioselectieve reactie}\index{regioselectieve reactie}; een
reactie die meerdere \omterm{def:b1:stereochemistry-in-depth:isomers}{stereo-isomeren} kan geven maar er vooral één vormt,
is een \emph{stereoselectieve reactie}\index{stereoselectieve reactie}.
\end{definition}

\begin{proposition}[Regel van Zaitsev]\label{prop:b1:elimination:zaitsev}
Wanneer meerdere $\beta$-waterstofatomen verwijderd kunnen worden, is het
hoofdalkeen meestal het meest gesubstitueerde (het stabielste), en eerder
het \cip{E}- dan het \cip{Z}-isomeer (de \emph{regel van
Zaitsev}\index{regel van Zaitsev}). Omvangrijke basen, en geometrische
beperkingen zoals de trans-diaxiale voorwaarde, kunnen haar opheffen.
\end{proposition}

\begin{proof}
Alkylgroepen op de koolstofatomen van een dubbele binding stabiliseren
die (zoals ze \omterm{def:b1:electronic-effects:intermediates}{carbokationen} stabiliseren), en \cip{E}-alkenen vermijden
het gedrang van \textit{cis}-groepen. De \omterm{def:b1:elementary-steps:energy-profile}{overgangstoestanden} van E2 en E1
hebben gedeeltelijk het karakter van een dubbele binding, zodat het
stabielere alkeen sneller ontstaat. Een omvangrijke base bereikt de
minder gehinderde waterstofatomen gemakkelijker; en een waterstofatoom dat
niet \omterm{def:b1:elimination:periplanar}{antiperiplanair} ten opzichte van de \omterm{def:b1:electronic-effects:nucleophile}{vertrekkende groep} kan staan,
wordt door E2 helemaal niet verwijderd, hoe stabiel het alkeen dat het
zou geven ook is.
\end{proof}

\begin{method}[Het alkeen voorspellen]\label{met:b1:elimination:predict-alkene}
\begin{enumerate}
  \item Som de $\beta$-koolstofatomen op die waterstofatomen dragen.
  \item Houd voor E2 alleen de waterstofatomen over die \omterm{def:b1:elimination:periplanar}{antiperiplanair}
        ten opzichte van Y kunnen staan (op een ring: trans-diaxiaal, in
        een \omterm{def:b1:stereochemistry-in-depth:conformations}{stoel} die werkelijk kan ontstaan).
  \item Geef daaronder de voorkeur aan het waterstofatoom dat het meest
        gesubstitueerde alkeen geeft (Zaitsev), tenzij de base omvangrijk
        is.
  \item Teken voor elk de \omterm{def:b1:stereochemistry-in-depth:representations}{newmanprojectie} in de \omterm{def:b1:stereochemistry-in-depth:conformations}{anticonformatie} om de
        \cip{E}/\cip{Z}-geometrie af te lezen.
\end{enumerate}
\end{method}

\section{Substitutie of eliminatie?}

\begin{proposition}[Substitutie tegenover eliminatie]\label{prop:b1:elimination:sn-vs-e}
Eliminatie wordt begunstigd door sterke, omvangrijke basen, door een hoge
temperatuur en door tertiaire of volle \omterm{def:b1:nucleophilic-substitution:substitution}{substraten}; substitutie door
goede \omterm{def:b1:electronic-effects:nucleophile}{nucleofielen} die \omterm{def:b1:predominant-reaction:strong-acid}{zwakke basen} zijn, door een lage temperatuur en
door primaire \omterm{def:b1:nucleophilic-substitution:substitution}{substraten}.
\end{proposition}

\begin{proof}
De sterkte van de base begunstigt de aanval op waterstof; omvang hindert
de aanval op een vol koolstofatoom, maar niet die op een vrijliggend
waterstofatoom. Een \omterm{def:b1:electronic-effects:carbon-class}{tertiair koolstofatoom} is onbereikbaar voor SN2, en
zijn \omterm{def:b1:electronic-effects:intermediates}{carbokation} verliest gemakkelijk protonen. Eliminatie maakt van één
molecuul er twee of drie (alkeen, \omterm{def:b1:electronic-effects:nucleophile}{vertrekkende groep}, geprotoneerde
base): ze wordt begunstigd naarmate de temperatuur stijgt, een resultaat
dat in het deel van jaar~2 verantwoord wordt.
\end{proof}

\begin{omfigure}
\begin{tabular}{llll}
\hline
\omterm{def:b1:nucleophilic-substitution:substitution}{substraat} & \omterm{def:b1:predominant-reaction:strong-acid}{sterke base}, klein & \omterm{def:b1:predominant-reaction:strong-acid}{sterke base}, omvangrijk & \omterm{def:b1:predominant-reaction:strong-acid}{zwakke base}, goed \omterm{def:b1:electronic-effects:nucleophile}{nucleofiel} \\
\hline
primair & SN2 (wat E2) & E2 & SN2 \\
secundair & E2 (en SN2) & E2 & SN2 (aprotisch), SN1/E1 (protisch) \\
tertiair & E2 & E2 & SN1/E1 (protisch) \\
\hline
\end{tabular}

\smallskip
{\small Een eerste oriëntatie: de hoofdreactie in elk geval. Verwarmen
begunstigt in elke kolom de eliminatie.}
\end{omfigure}

\section{Oefeningen}

\begin{exercise}[$\star$]\label{exo:b1:elimination:1}
Geef de alkenen die kunnen ontstaan door eliminatie van HBr uit
2-broombutaan, uit 2-broom-2-methylbutaan en uit broomcyclohexaan, en het
alkeen dat volgens de \omterm{prop:b1:elimination:zaitsev}{regel van Zaitsev} het hoofdproduct is.
\end{exercise}

\begin{exercise}[$\star$]\label{exo:b1:elimination:2}
Schrijf de \omterm{def:b1:rate-laws:order}{snelheidswetten} van E1 en E2 op. Hoe zou je ze experimenteel
van elkaar onderscheiden voor 2-broom-2-methylpropaan in ethanol met
ethoxide?
\end{exercise}

\begin{exercise}[$\star$]\label{exo:b1:elimination:3}
Teken de \omterm{def:b1:stereochemistry-in-depth:representations}{newmanprojecties} van 2-broombutaan langs C2--C3 in zijn drie
\omterm{def:b1:stereochemistry-in-depth:conformations}{gestaffelde conformaties}, en zeg welke E2 kunnen ondergaan.
\end{exercise}

\begin{exercise}[$\star$]\label{exo:b1:elimination:4}
Voorspel de hoofdreactie (SN1, SN2, E1, E2): broomethaan met
natriumhydroxide in water bij \qty{25}{\celsius}; 2-broom-2-methylpropaan
met kalium-tert-butoxide; 2-broompropaan met natriumjodide in propanon;
2-broom-2-methylpropaan in hete ethanol.
\end{exercise}

\begin{exercise}[$\star\star$]\label{exo:b1:elimination:5}
Toon met \omterm{def:b1:stereochemistry-in-depth:representations}{newmanprojecties} aan dat de E2-eliminatie van HBr uit de twee
\omterm{def:b1:stereochemistry-in-depth:enantiomers}{diastereomeren} van 2-broom-3-methylpentaan die het waterstofatoom van C3
kunnen verliezen, de twee \omterm{def:b1:stereochemistry-in-depth:isomers}{stereo-isomeren} van 3-methylpent-2-een geeft.
\end{exercise}

\begin{exercise}[$\star\star$]\label{exo:b1:elimination:6}
Schrijf het mechanisme op van de dehydratatie van 2-methylbutaan-2-ol in
heet geconcentreerd zwavelzuur (E1) en geef het hoofdalkeen.
\end{exercise}

\begin{exercise}[$\star\star$]\label{exo:b1:elimination:7}
Leg uit waarom \textit{trans}-1-broom-4-tert-butylcyclohexaan zeer traag
volgens E2 reageert, terwijl zijn \textit{cis}-isomeer snel reageert. (De
tert-butylgroep blijft \omterm{def:b1:stereochemistry-in-depth:conformations}{equatoriaal}.)
\end{exercise}

\begin{exercise}[$\star\star$]\label{exo:b1:elimination:8}
2-Broom-2-methylbutaan geeft met natriumethoxide vooral
2-methylbut-2-een; met kalium-tert-butoxide vooral 2-methylbut-1-een.
Verklaar.
\end{exercise}

\begin{exercise}[$\star\star$]\label{exo:b1:elimination:9}
Waarom wordt er bij broommethaan nooit eliminatie waargenomen? En bij
1-broom-2,2-dimethylpropaan via E2?
\end{exercise}

\begin{exercise}[$\star\star\star$]\label{exo:b1:elimination:10}
In ethanol bij \qty{25}{\celsius} geeft 2-broom-2-methylpropaan een
mengsel van de ether (SN1) en het alkeen (E1), waarvan de verhoudingen
niet afhangen van de concentratie van het \omterm{def:b1:nucleophilic-substitution:substitution}{substraat}. Toon aan dat beide
producten uit hetzelfde \omterm{def:b1:electronic-effects:intermediates}{carbokation} komen en dat hun verhouding de
verhouding van twee \omterm{def:b1:rate-laws:order}{snelheidsconstanten} is.
\end{exercise}

\begin{exercise}[$\star\star\star$]\label{exo:b1:elimination:11}
Eén \omterm{def:b1:stereochemistry-in-depth:enantiomers}{diastereomeer} van 1,2-dibroom-1,2-difenylethaan (de mesovorm) geeft
door E2 met ethoxide één enkel \omterm{def:b1:stereochemistry-in-depth:isomers}{stereo-isomeer} van
1-broom-1,2-difenyletheen. Bepaal welk, met een \omterm{def:b1:stereochemistry-in-depth:representations}{newmanprojectie}.
\end{exercise}

\begin{exercise}[$\star\star\star$]\label{exo:b1:elimination:12}
Leg met de energieën van \cref{ch:b1:stereochemistry-in-depth} uit
waarom een E2-reactie op een cyclohexaan honderden keren vertraagd kan
worden wanneer de \omterm{def:b1:electronic-effects:nucleophile}{vertrekkende groep} in een volle \omterm{def:b1:stereochemistry-in-depth:conformations}{stoel} \omterm{def:b1:stereochemistry-in-depth:conformations}{axiaal} moet
worden. Druk de snelheid uit als het product van een fractie conformeren
en een \omterm{def:b1:rate-laws:order}{snelheidsconstante}.
\end{exercise}

\section{Probleem: menthyl- en neomenthylchloride}

\begin{problem}[{Weekendprobleem --- de stoelen van twee diastereomere
chloriden, de antiperiplanaire waterstofatomen die in elk beschikbaar
zijn, de gevormde alkenen, en de verhouding van hun E2-snelheden}]
\label{pb:b1:elimination:1}
Menthylchloride en neomenthylchloride zijn de chloriden van menthol en
neomenthol (\cref{ch:b1:stereochemistry-in-depth}): op de ring draagt C1
het Cl, C2 de isopropylgroep, C5 de methylgroep. In menthylchloride
kunnen de drie groepen allemaal \omterm{def:b1:stereochemistry-in-depth:conformations}{equatoriaal} staan; in neomenthylchloride
staat Cl \omterm{def:b1:stereochemistry-in-depth:conformations}{axiaal} wanneer de twee alkylgroepen \omterm{def:b1:stereochemistry-in-depth:conformations}{equatoriaal} staan. Beide
worden behandeld met natriumethoxide in ethanol, een E2-reactie. Gegevens
van het probleem: \qty{95}{\percent} van het neomenthylchloride zit in de
\omterm{def:b1:stereochemistry-in-depth:conformations}{stoel} met Cl \omterm{def:b1:stereochemistry-in-depth:conformations}{axiaal}; voor menthylchloride vormt de \omterm{def:b1:stereochemistry-in-depth:conformations}{stoel} met Cl \omterm{def:b1:stereochemistry-in-depth:conformations}{axiaal}
(alle drie de groepen \omterm{def:b1:stereochemistry-in-depth:conformations}{axiaal}) een fractie \num{5.0e-3}; neem aan dat elk
\omterm{def:b1:elimination:periplanar}{antiperiplanair} waterstofatoom van een reactieve \omterm{def:b1:stereochemistry-in-depth:conformations}{stoel} met dezelfde
\omterm{def:b1:rate-laws:order}{snelheidsconstante} reageert. Neomenthylchloride geeft
\qty{78}{\percent} van het trigesubstitueerde alkeen.

\medskip
\noindent\textbf{Deel I --- De \omterm{def:b1:stereochemistry-in-depth:conformations}{stoelen}.}
\begin{enumerate}
  \item Teken de voorkeursstoel van menthylchloride.
  \item Teken de voorkeursstoel van neomenthylchloride.
  \item Zijn de twee chloriden \omterm{def:b1:stereochemistry-in-depth:enantiomers}{enantiomeren} of \omterm{def:b1:stereochemistry-in-depth:enantiomers}{diastereomeren}?
  \item Formuleer de geometrische voorwaarde van E2 op een ring.
  \item In welke \omterm{def:b1:stereochemistry-in-depth:conformations}{stoel} kan menthylchloride reageren? Teken die.
  \item Waarom is die \omterm{def:b1:stereochemistry-in-depth:conformations}{stoel} zeldzaam? Gebruik de energiekosten van \omterm{def:b1:stereochemistry-in-depth:conformations}{axiale}
        groepen.
\end{enumerate}

\medskip
\noindent\textbf{Deel II --- De anti-waterstofatomen.}
\begin{enumerate}[resume]
  \item Som de $\beta$-koolstofatomen van het chloride op (C2 en C6) en
        hun waterstofatomen.
  \item Welke waterstofatomen staan in de reactieve \omterm{def:b1:stereochemistry-in-depth:conformations}{stoel} van
        neomenthylchloride \omterm{def:b1:elimination:periplanar}{antiperiplanair} ten opzichte van Cl?
  \item Staat het waterstofatoom van C2 in de reactieve \omterm{def:b1:stereochemistry-in-depth:conformations}{stoel} van
        menthylchloride \omterm{def:b1:stereochemistry-in-depth:conformations}{axiaal}? Staat er een waterstofatoom van C6
        \omterm{def:b1:stereochemistry-in-depth:conformations}{axiaal}?
  \item Hoeveel \omterm{def:b1:elimination:periplanar}{antiperiplanaire} waterstofatomen biedt elk chloride?
  \item Teken een \omterm{def:b1:stereochemistry-in-depth:representations}{newmanprojectie} langs C1--C6 voor de reactieve \omterm{def:b1:stereochemistry-in-depth:conformations}{stoel} van
        menthylchloride.
  \item Waarom zou een \omterm{def:b1:elimination:periplanar}{synperiplanaire} eliminatie de reactie niet kunnen
        redden?
\end{enumerate}

\medskip
\noindent\textbf{Deel III --- De producten.}
\begin{enumerate}[resume]
  \item Benoem het alkeen dat ontstaat door het waterstofatoom van C2 te
        verwijderen, en het alkeen dat ontstaat door een waterstofatoom
        van C6 te verwijderen.
  \item Welk ervan begunstigt de \omterm{prop:b1:elimination:zaitsev}{regel van Zaitsev}, en waarom?
  \item Wat geeft neomenthylchloride? Is de verhouding 78/22 in
        overeenstemming met de \omterm{prop:b1:elimination:zaitsev}{regel van Zaitsev}?
  \item Wat geeft menthylchloride? Is dat in overeenstemming met de \omterm{prop:b1:elimination:zaitsev}{regel
        van Zaitsev}?
  \item Welke reactie is stereospecifiek, regioselectief, of beide?
  \item Waarom is substitutie (SN2) met ethoxide hier ondergeschikt?
\end{enumerate}

\medskip
\noindent\textbf{Deel IV --- De snelheden.}
\begin{enumerate}[resume]
  \item Schrijf de snelheid van elke reactie op als functie van de fractie
        reactieve \omterm{def:b1:stereochemistry-in-depth:conformations}{stoel} en het aantal \omterm{def:b1:elimination:periplanar}{antiperiplanaire} waterstofatomen.
  \item Bereken de relatieve snelheid van neomenthylchloride.
  \item Bereken de relatieve snelheid van menthylchloride.
  \item Bereken de verhouding van de snelheden
        $k(\text{neomenthyl})/k(\text{menthyl})$.
\end{enumerate}
\end{problem}
